% Aula 03 --- demonstrar o atalho do LOOCV para ajustes lineares.
O LOOCV (\emph{Leave-One-Out Cross Validation}) parece caro: $n$ ajustes, um para cada observação deixada de fora.
Já comentamos que, para ajustes lineares, um único ajuste basta --- mas não demonstramos.

Vamos demonstrar no caso da regressão linear por mínimos quadrados. A matriz de
delineamento $\mathbb{X}$ tem linha $i$ igual a $\X_i^\top$, e o ajuste com
\emph{todos} os dados é $\widehat g(\x)=\x^\top\widehat\bbeta$, com
$\widehat\bbeta=(\mathbb{X}^\top\mathbb{X})^{-1}\mathbb{X}^\top\bm y$. Chame de
$\widehat g_{-i}$ o ajuste feito sem a $i$-ésima observação. A fórmula a demonstrar
é
\[
  \riscohat_{\text{LOO}}
   = \frac1n\sum_{i=1}^{n}\left(\frac{Y_i-\widehat g(\X_i)}{1-h_{ii}}\right)^{2},
\]
em que $h_{ii}$ é o $i$-ésimo elemento da diagonal da matriz
$\bm H=\mathbb{X}(\mathbb{X}^\top\mathbb{X})^{-1}\mathbb{X}^\top$.

Suponha que $\mathbb{X}$ tenha posto cheio, e que continue tendo quando se tira
qualquer uma de suas linhas. Assim todos os ajustes são únicos, e $h_{ii}\neq 1$:
dividir por $1-h_{ii}$ nunca é problema.

A demonstração vem em etapas. O item (a) mostra o que são os números $h_{i\ell}$;
os itens (b) e (c) usam um truque para escrever $\widehat g_{-i}(\X_i)$ em termos de
$\widehat g$; o (d) faz a conta; e o (e) interpreta o resultado.
\begin{itens}
  \item \pts{0,6} \emph{De onde vêm os pesos.} Mostre que o vetor de valores
        ajustados, $\widehat{\bm y}=\big(\widehat g(\X_1),\dots,\widehat g(\X_n)\big)^\top$,
        é igual a $\bm H\bm y$, e conclua, olhando a coordenada $i$, que
        \[
          \widehat g(\X_i)=\sum_{\ell=1}^{n} h_{i\ell}\,Y_\ell ,
        \]
        em que $h_{i\ell}$ é a entrada $(i,\ell)$ de $\bm H$. Dica: como a linha
        $i$ de $\mathbb{X}$ é $\X_i^\top$, vale $\widehat{\bm y}=\mathbb{X}\widehat\bbeta$.

        Em palavras: a predição no ponto $\X_i$ é uma combinação de todas as
        respostas, inclusive da própria $Y_i$, que entra com peso $h_{ii}$. De que
        esses pesos dependem? E de que eles \emph{não} dependem?
  \item \pts{1,0} \emph{O truque.} Imagine o conjunto de dados em que a resposta
        $Y_i$ foi trocada por $\widetilde Y_i=\widehat g_{-i}(\X_i)$ --- o valor que o
        ajuste sem a observação $i$ prevê ali ---, com as covariáveis e as demais
        respostas intactas. Mostre que o ajuste de mínimos quadrados nesse conjunto
        modificado é exatamente $\widehat g_{-i}$.

        Antes do caso geral, veja o truque funcionar no modelo só com intercepto,
        em que o ajuste é a média: troque $Y_i$ pela média das outras $n-1$
        respostas e confira que a média das $n$ passa a ser justamente a média das
        outras. Para o caso geral, um caminho: escreva o RSS do conjunto modificado
        como a soma das $n-1$ parcelas de sempre com a parcela da observação $i$;
        veja o que $\widehat g_{-i}$ faz com cada uma das duas partes; e conclua
        que nenhuma função linear consegue um RSS menor, o que basta, porque o
        ajuste é único.
  \item \pts{0,8} \emph{A mesma matriz, duas vezes.} O conjunto modificado tem as
        mesmas covariáveis que o original, e portanto a mesma matriz $\bm H$.
        Aplique (a) a ele e use (b) para mostrar que
        \[
          \widehat g_{-i}(\X_i)=\sum_{\ell\neq i} h_{i\ell}\,Y_\ell + h_{ii}\,\widetilde Y_i .
        \]
        Depois some e subtraia $h_{ii}Y_i$ para chegar a
        \[
          \widehat g_{-i}(\X_i) = \widehat g(\X_i) + h_{ii}\big(\widetilde Y_i - Y_i\big).
        \]
  \item \pts{0,7} \emph{A conta.} Na igualdade de (c), $\widetilde Y_i$ aparece dos
        dois lados, porque o lado esquerdo é o próprio $\widetilde Y_i$. Subtraia
        $Y_i$ dos dois lados, junte os termos com $\widetilde Y_i-Y_i$ e isole o
        resíduo de validação $Y_i-\widehat g_{-i}(\X_i)$. Conclua a fórmula do
        enunciado.
  \item \pts{0,9} \emph{Lendo o resultado.} A matriz $\bm H$ tem duas propriedades.
        É simétrica, porque $\mathbb{X}^\top\mathbb{X}$ é simétrica e sua inversa
        também. E vale $\bm H\bm H=\bm H$, o que se confere multiplicando:
        \[
          \bm H\bm H
           = \mathbb{X}(\mathbb{X}^\top\mathbb{X})^{-1}
             \underbrace{\mathbb{X}^\top\mathbb{X}(\mathbb{X}^\top\mathbb{X})^{-1}}_{=\;\mathbb{I}}
             \mathbb{X}^\top
           = \bm H .
        \]
        Em palavras: $\bm H$ projeta $\bm y$ no espaço gerado pelas colunas de
        $\mathbb{X}$, e projetar de novo o que já foi projetado não muda nada.
        Usando as duas propriedades, a entrada $(i,i)$ de $\bm H\bm H=\bm H$ dá
        \[
          h_{ii}=\sum_{\ell=1}^{n} h_{i\ell}^2 = h_{ii}^2+\sum_{\ell\neq i} h_{i\ell}^2 .
        \]
        Conclua que $0\le h_{ii}\le 1$ e que, portanto, o resíduo de validação
        nunca é menor, em módulo, que o de treino. O que isso confirma sobre o erro
        de treino como estimativa do risco? E o que significa uma observação com
        $h_{ii}$ próximo de $1$?
\end{itens}

\begin{solucao}
\textbf{(a)} Como a linha $i$ de $\mathbb{X}$ é $\X_i^\top$, a coordenada $i$ de
$\mathbb{X}\widehat\bbeta$ é $\X_i^\top\widehat\bbeta=\widehat g(\X_i)$; logo
$\widehat{\bm y}=\mathbb{X}\widehat\bbeta$ e, substituindo $\widehat\bbeta$,
\[
  \widehat{\bm y}
  =\underbrace{\mathbb{X}(\mathbb{X}^\top\mathbb{X})^{-1}\mathbb{X}^\top}_{=\;\bm H}\,\bm y
  =\bm H\bm y .
\]
A coordenada $i$ de $\bm H\bm y$ é a linha $i$ de $\bm H$ vezes $\bm y$; igualando as
coordenadas $i$ dos dois lados, $\widehat g(\X_i)=\sum_\ell h_{i\ell}Y_\ell$.

Os pesos são entradas de $\bm H$, que é construída só com $\mathbb{X}$: eles
\textbf{dependem das covariáveis e não dependem das respostas}. Troque o vetor $\bm y$
mantendo as covariáveis, e a predição em $\X_i$ continua sendo a mesma combinação,
agora das respostas novas. É essa propriedade --- o ajuste é uma função \emph{linear}
de $\bm y$ --- que o item (c) usa. (No modelo só com intercepto, $\mathbb{X}=\bm1$, todas
as entradas de $\bm H$ valem $1/n$, e a combinação é a média.)

\smallskip
\textbf{(b)} \emph{O caso da média.} Seja $\bar Y_{-i}$ a média das outras $n-1$
respostas. Trocando $Y_i$ por $\widetilde Y_i=\bar Y_{-i}$, a soma das $n$ respostas vira
$(n-1)\bar Y_{-i}+\bar Y_{-i}=n\bar Y_{-i}$, e a nova média é $\bar Y_{-i}$. A observação
trocada foi posta exatamente onde o ajuste sem ela já estava, e por isso não puxa o
ajuste para lado nenhum.

\emph{O caso geral} é o mesmo argumento, escrito com o RSS. Para uma função linear
$g(\x)=\x^\top\bbeta$ qualquer, o RSS do conjunto modificado se parte em duas:
\[
  \widetilde{\mathrm{RSS}}(g)
  =\underbrace{\sum_{\ell\neq i}\big(Y_\ell-g(\X_\ell)\big)^2}_{A(g)}
  +\underbrace{\big(\widetilde Y_i-g(\X_i)\big)^2}_{B(g)\;\ge\;0} .
\]
Por definição, $\widehat g_{-i}$ minimiza $A$ entre as funções lineares --- e está bem
definido, porque $\mathbb{X}$ sem a linha $i$ tem posto cheio. E $B(\widehat g_{-i})=0$,
porque $\widetilde Y_i$ foi escolhido como $\widehat g_{-i}(\X_i)$. Logo, para toda $g$
linear,
\[
  \widetilde{\mathrm{RSS}}(g)=A(g)+B(g)\;\ge\;A(\widehat g_{-i})+0
  =\widetilde{\mathrm{RSS}}(\widehat g_{-i}):
\]
$\widehat g_{-i}$ minimiza o RSS modificado. Como o conjunto modificado tem a mesma
matriz $\mathbb{X}$, de posto cheio, esse minimizador é único; é, portanto, o ajuste de
mínimos quadrados do conjunto modificado.

\smallskip
\textbf{(c)} Aplique (a) ao conjunto modificado. As covariáveis são as mesmas, então
$\bm H$ é a mesma, e a predição em $\X_i$ é a mesma combinação de antes, com
$\widetilde Y_i$ no lugar de $Y_i$. Por (b), essa predição é $\widehat g_{-i}(\X_i)$:
\[
  \widehat g_{-i}(\X_i)=\sum_{\ell\neq i}h_{i\ell}Y_\ell+h_{ii}\widetilde Y_i .
\]
Somando e subtraindo $h_{ii}Y_i$, a soma se completa e aparece $\widehat g(\X_i)$:
\[
  \widehat g_{-i}(\X_i)
  =\underbrace{\sum_{\ell}h_{i\ell}Y_\ell}_{=\;\widehat g(\X_i)}+h_{ii}\big(\widetilde Y_i-Y_i\big).
\]

\smallskip
\textbf{(d)} O lado esquerdo de (c) é o próprio $\widetilde Y_i$. Subtraindo $Y_i$ dos
dois lados, a diferença $\widetilde Y_i-Y_i$ aparece nos dois:
\[
  \widetilde Y_i-Y_i=\big(\widehat g(\X_i)-Y_i\big)+h_{ii}\big(\widetilde Y_i-Y_i\big)
  \iff
  (1-h_{ii})\big(\widetilde Y_i-Y_i\big)=\widehat g(\X_i)-Y_i .
\]
Como $h_{ii}\neq1$ (o item (e) confere), podemos dividir; trocando os sinais,
\[
  \boxed{\;Y_i-\widehat g_{-i}(\X_i)=\frac{Y_i-\widehat g(\X_i)}{1-h_{ii}}\;}
\]
Elevando ao quadrado e tirando a média em $i$, sai a fórmula do enunciado. O ganho está
no lado direito: só entra $\widehat g$, ajustado \emph{uma} vez, com todos os dados. Os
$n$ ajustes do LOOCV viraram um.

\smallskip
\textbf{(e)} Pela simetria, a entrada $(i,i)$ de $\bm H\bm H$ é
$\sum_\ell h_{i\ell}h_{\ell i}=\sum_\ell h_{i\ell}^2$; pela idempotência, ela vale
$h_{ii}$. Daí saem as duas desigualdades:
\[
  h_{ii}=\sum_\ell h_{i\ell}^2\;\ge\;0,
  \qquad
  h_{ii}=h_{ii}^2+\sum_{\ell\neq i}h_{i\ell}^2\;\ge\;h_{ii}^2
  \;\Longrightarrow\;h_{ii}(1-h_{ii})\ge0
  \;\Longrightarrow\;h_{ii}\le1 .
\]

\emph{E a hipótese exclui $h_{ii}=1$.} Suponha $h_{ii}=\X_i^\top(\mathbb{X}^\top\mathbb{X})^{-1}\X_i=1$
e tome $\bm v=(\mathbb{X}^\top\mathbb{X})^{-1}\X_i$, que não é nulo (senão $h_{ii}$ seria
$0$). Sem a linha $i$, a matriz de delineamento $\mathbb{X}_{-i}$ satisfaz
$\mathbb{X}_{-i}^\top\mathbb{X}_{-i}=\mathbb{X}^\top\mathbb{X}-\X_i\X_i^\top$, e então
\[
  \mathbb{X}_{-i}^\top\mathbb{X}_{-i}\,\bm v=\X_i-\X_i\,(\X_i^\top\bm v)=\X_i-\X_i\,h_{ii}=\bm0 :
\]
$\mathbb{X}_{-i}$ perderia o posto, contra a hipótese. (A recíproca também vale, e o
exemplo típico é uma coluna indicadora que só a observação $i$ tem: sem ela, a coluna
zera.)

Com $0\le h_{ii}<1$, o fator $1/(1-h_{ii})$ é pelo menos $1$, e (d) dá
\[
  \abs{Y_i-\widehat g_{-i}(\X_i)}\;\ge\;\abs{Y_i-\widehat g(\X_i)}
  \qquad\text{para toda observação } i .
\]
Tirando a média dos quadrados, $\riscohat_{\text{LOO}}$ nunca fica abaixo do erro
quadrático médio de treino. Isso confirma, de um jeito bem concreto, que o erro de
treino é \textbf{otimista} --- e não só em média: em toda amostra, observação por
observação, e com o tamanho da correção escrito.

O $h_{ii}$ mede a \textbf{alavancagem}: o peso da própria $Y_i$ na predição feita em
$\X_i$. Perto de $1$, essa predição é quase só $Y_i$ --- o ajuste passa colado no
ponto, e o resíduo de treino sai pequeno e engana. O fator $1/(1-h_{ii})$ fica grande
justamente para desfazer o engano: sem a observação $i$, as outras quase não informam
o que acontece em $\X_i$. Para ter uma escala,
$\sum_ih_{ii}=\operatorname{tr}(\bm H)=\operatorname{tr}\big((\mathbb{X}^\top\mathbb{X})^{-1}\mathbb{X}^\top\mathbb{X}\big)$
é o número de colunas de $\mathbb{X}$; a alavancagem \emph{média} é, portanto, o número
de parâmetros dividido por $n$, e "alta" quer dizer muito acima disso.

Um último comentário, de método: repare onde cada hipótese trabalhou. A linearidade do
ajuste nas respostas entrou em (a) e (c); o fato de $\widehat g_{-i}$ ser um ajuste de
mínimos quadrados \emph{na mesma classe de funções}, em (b); o posto cheio, na
unicidade de (b) e em $h_{ii}\neq1$; a simetria e a idempotência de $\bm H$, só em (e).
Um método que não tenha alguma dessas propriedades não herda a fórmula
automaticamente.
\end{solucao}
